% ==============================================================================
% PHỤ LỤC E: TUYỂN TẬP BÀI TẬP LỚN & ĐỒ ÁN ỨNG DỤNG AI TỪ CON SỐ 0
% ==============================================================================
% Phong cách: Sách chuyên khảo kinh điển (Classic Academic & Engineering Monograph)
% Tiêu chuẩn: Từ số 0, phân tích bản chất toán học, sơ đồ kiến trúc, mã nguồn chuẩn
% ==============================================================================

\chapter{Tuyển Tập Bài Tập Lớn \& Đồ Án Ứng Dụng AI Toàn Diện (Capstone Workbook)}
\label{ch:appendix-workbook}

Để thực sự làm chủ năng lực của một Kỹ sư AI cao cấp, việc nắm vững lý thuyết trừu tượng là chưa đủ. Bạn cần trực tiếp nhúng tay vào việc giải quyết các bài toán kỹ nghệ thế giới thực từ con số 0: Từ việc đọc từng gói tin mạng thô, giải tích các phương trình kiểm toán, đến việc thiết kế luồng suy luận của AI Agent có khả năng tự sửa lỗi và tích hợp hệ thống phòng thủ đa tầng.

Phụ lục chuyên khảo này thiết kế 6 đồ án ứng dụng AI cấp độ doanh nghiệp hoàn chỉnh. Mỗi đồ án đều được xây dựng theo quy trình kỹ nghệ 5 bước chuẩn mực:
\begin{enumerate}
  \item \textbf{Khởi điểm từ con số 0}: Định nghĩa các khái niệm nền tảng về mạng, tài chính, y tế hoặc thuật toán mà không giả định người đọc đã biết trước.
  \item \textbf{Phân tích động cơ \& Mô hình toán học}: Giải thích cặn kẽ tại sao phải áp dụng công thức giải tích hoặc xác suất cụ thể.
  \item \textbf{Kiến trúc hệ thống \& Luồng dữ liệu}: Sơ đồ luồng xử lý và tương tác giữa các microservices.
  \item \textbf{Triển khai mã nguồn Python chuẩn mực}: Code hoàn chỉnh có Type Hints, Logging, xử lý ngoại lệ và sẵn sàng kiểm thử.
  \item \textbf{Bộ câu hỏi phản biện bảo vệ đồ án}: Rèn luyện tư duy thiết kế hệ thống và đánh đổi kỹ thuật.
\end{enumerate}

---

\section{Đồ Án 1: Trung Tâm Điều Hành An Ninh Mạng AI Tự Hành (Autonomous AI-SOC)}
\label{sec:workbook-ai-soc}

\subsection{Khởi Điểm Từ Con Số 0: Mạng Máy Tính \& Cơ Chế Tấn Công}
Mọi hoạt động trên Internet đều dựa trên việc truyền nhận các gói tin (\textit{Packets}) qua giao thức TCP/IP. Mỗi gói tin chứa thông tin địa chỉ IP nguồn (\textit{Source IP}), IP đích (\textit{Destination IP}), Cổng nguồn/đích (\textit{Port}), và cờ điều khiển giao vận (\textit{SYN, ACK, FIN, RST}).

Khi kẻ tấn công muốn làm tê liệt một máy chủ AI Serving, hình thức phổ biến nhất là tấn công từ chối dịch vụ phân tán (DDoS qua SYN Flood): Kẻ tấn công gửi hàng triệu gói tin TCP SYN từ nhiều địa chỉ IP giả mạo nhưng không bao giờ gửi lại gói tin ACK cuối cùng để hoàn tất quy trình bắt tay ba bước (3-way handshake). Hàng đợi kết nối (\textit{Backlog Queue}) của hệ điều hành Linux lập tức bị tràn ngập, khiến các người dùng hợp lệ không thể kết nối vào hệ thống.

\subsection{Mô Hình Toán Học: Entropy Phân Phối Gói Tin Thời Gian Thực}
Để phát hiện tấn công mạng tự động mà không bị phụ thuộc vào các tập luật tĩnh cố định, ta sử dụng lý thuyết thông tin Shannon để đo lường độ hỗn loạn (\textit{Entropy}) của phân phối địa chỉ IP trong một cửa sổ thời gian trượt $\Delta t = 10\text{ giây}$.

Giả sử trong cửa sổ $\Delta t$, hệ thống ghi nhận $M$ gói tin tới từ $K$ địa chỉ IP nguồn khác nhau. Tần suất xuất hiện của địa chỉ IP thứ $i$ là:
\begin{equation}
p_i = \frac{m_i}{M}, \quad \text{với } \sum_{i=1}^K p_i = 1
\end{equation}
Entropy Shannon của phân phối địa chỉ IP nguồn:
\begin{equation}
H(X) = - \sum_{i=1}^K p_i \log_2(p_i)
\end{equation}

\textbf{Quy luật toán học}:
\begin{itemize}
  \item Ở trạng thái hoạt động bình thường, các gói tin đến từ hàng ngàn người dùng ngẫu nhiên phân tán, phân phối $p_i$ tương đối đồng đều $\implies H(X)$ đạt giá trị cao (tiệm cận $\log_2 K$).
  \item Khi xảy ra tấn công DoS từ một botnet tập trung, một lượng lớn gói tin dồn dập xuất phát từ một nhóm nhỏ IP hoặc tập trung đánh vào duy nhất một cổng $\implies$ Phân phối bị méo mó nghiêm trọng $\implies H(X)$ sụt giảm đột ngột xuống dưới ngưỡng cảnh báo $H_{\text{threshold}}$.
\end{itemize}

\subsection{Triển Khai Mã Nguồn Python: Bộ Phát Hiện Anomaly \& Agent Cách Ly IP}

\begin{lstlisting}[language=python]
import math
import time
from collections import Counter
from dataclasses import dataclass

@dataclass
class NetworkPacket:
    timestamp: float
    src_ip: str
    dst_ip: str
    dst_port: int
    protocol: str
    tcp_flags: str

class AutonomousSOCDetector:
    def __init__(self, entropy_threshold: float = 1.5, window_seconds: float = 10.0):
        self.entropy_threshold = entropy_threshold
        self.window_seconds = window_seconds
        self.packet_buffer: list[NetworkPacket] = []
        self.blocked_ips: set[str] = set()

    def ingest_packet(self, packet: NetworkPacket) -> None:
        """Thu thap goi tin va loai bo cac goi tin nam ngoai cua so thoi gian."""
        self.packet_buffer.append(packet)
        cutoff = packet.timestamp - self.window_seconds
        self.packet_buffer = [p for p in self.packet_buffer if p.timestamp >= cutoff]

    def compute_source_entropy(self) -> float:
        """Tinh toan Entropy Shannon cua cac IP nguon trong cua so hien tai."""
        if not self.packet_buffer:
            return 0.0
        total_packets = len(self.packet_buffer)
        counts = Counter(p.src_ip for p in self.packet_buffer)
        entropy = 0.0
        for count in counts.values():
            p_i = count / total_packets
            if p_i > 0:
                entropy -= p_i * math.log2(p_i)
        return entropy

    def evaluate_and_respond(self) -> dict:
        """Danh gia trang thai mang va tu dong kich hoat luat tuong lua iptables."""
        entropy = self.compute_source_entropy()
        total_packets = len(self.packet_buffer)
        
        # Phat hien bat thuong khi luu luong lon va entropy sup giam dot ngot
        if total_packets > 50 and entropy < self.entropy_threshold:
            # Tim IP co mat do tan cong cao nhat
            counts = Counter(p.src_ip for p in self.packet_buffer)
            top_attacker, top_count = counts.most_common(1)[0]
            
            # Kiem tra xem da chan chua
            if top_attacker not in self.blocked_ips:
                self.blocked_ips.add(top_attacker)
                iptables_rule = f"iptables -I INPUT -s {top_attacker} -j DROP"
                return {
                    "alert": True,
                    "status": "ATTACK_DETECTED",
                    "entropy": entropy,
                    "attacker_ip": top_attacker,
                    "mitigation_action": iptables_rule
                }
        return {"alert": False, "status": "NORMAL", "entropy": entropy}
\end{lstlisting}

\subsection{Bộ Câu Hỏi Phản Biện Bảo Vệ Đồ Án}
\begin{enumerate}
  \item \textbf{Câu hỏi}: Nếu kẻ tấn công sử dụng kỹ thuật IP Spoofing tạo ra hàng triệu IP nguồn ngẫu nhiên thì entropy sẽ tăng cao thay vì giảm, giải pháp Entropy Shannon sẽ xử lý thế nào?
  \item \textbf{Trả lời}: Khi IP nguồn bị giả mạo ngẫu nhiên (phân tán cực đại), ta chuyển sang tính Entropy của \textit{Cổng đích (Destination Port)} hoặc tỷ lệ cờ TCP (\textit{SYN/ACK ratio}). Trong tấn công SYN Flood, tỷ lệ SYN không có ACK đi kèm sẽ chạm ngưỡng bất thường, giúp hệ thống phát hiện chính xác.
\end{enumerate}

---

\section{Đồ Án 2: Hệ Thống Đối Soát \& Kiểm Toán Báo Cáo Tài Chính Đa Phương Thức}
\label{sec:workbook-multimodal-auditor}

\subsection{Khởi Điểm Từ Con Số 0: Bản Chất Bảng Cân Đối Kế Toán}
Trong tài chính doanh nghiệp, Bảng Cân Đối Kế Toán (\textit{Balance Sheet}) phản ánh bức tranh sức khỏe tài chính tại một thời điểm cụ thể. Nguyên tắc bất di bất dịch của kế toán kép là Phương trình Cân đối Kế toán:
\begin{equation}
\text{Tổng Tài Sản (Total Assets)} = \text{Tổng Nợ Phải Trả (Liabilities)} + \text{Vốn Chủ Sở Hữu (Equity)}
\end{equation}
Khi một doanh nghiệp nộp báo cáo quét dạng PDF, các lỗi quang học OCR có thể nhận diện nhầm số $8$ thành số $3$ hoặc $0$, khiến phương trình kế toán bị lệch. Nhiệm vụ của Kỹ sư AI là xây dựng một pipeline trích xuất thị giác (Vision LLM) kết hợp một engine xác thực toán học có khả năng tự động truy vết sai lệch.

\subsection{Mô Hình Toán Học: Sai Số Cân Đối Kế Toán \& Khoảng Cách Levenshtein}
Hàm mất mát sai lệch cân đối kế toán:
\begin{equation}
\Delta_{\text{Balance}} = \left| A - (L + E) \right|
\end{equation}
Hệ thống chỉ phê duyệt báo cáo tài chính là toàn vẹn nếu sai số tương đối không vượt quá ngưỡng sai số số học:
\begin{equation}
\frac{\Delta_{\text{Balance}}}{A} \le 10^{-4} \quad (0.01\%)
\end{equation}

Đối với tên nhà cung cấp và số tài khoản ngân hàng, hệ thống áp dụng khoảng cách Levenshtein $\text{Lev}(s_1, s_2)$ để so khớp mờ với cơ sở dữ liệu danh mục được phê duyệt:
\begin{equation}
\text{Similarity}(s_1, s_2) = 1 - \frac{\text{Lev}(s_1, s_2)}{\max(|s_1|, |s_2|)}
\end{equation}

\subsection{Triển Khai Mã Nguồn Python: Engine Kiểm Toán Tài Chính}

\begin{lstlisting}[language=python]
from dataclasses import dataclass
from typing import Optional

@dataclass
class FinancialStatement:
    company_name: str
    total_assets: float
    total_liabilities: float
    equity: float
    reported_net_income: float

class FinancialAuditEngine:
    def __init__(self, tolerance_ratio: float = 0.0001):
        self.tolerance_ratio = tolerance_ratio

    def audit_balance_sheet(self, stmt: FinancialStatement) -> dict:
        """Kiem toan tinh can doi cua bao cao tai chinh theo phuong trinh co ban."""
        expected_assets = stmt.total_liabilities + stmt.equity
        discrepancy = abs(stmt.total_assets - expected_assets)
        
        # Tinh ty le sai lech tren tong tai san
        error_ratio = discrepancy / (stmt.total_assets + 1e-9)
        is_balanced = error_ratio <= self.tolerance_ratio
        
        audit_verdict = {
            "company": stmt.company_name,
            "total_assets": stmt.total_assets,
            "calculated_liabilities_equity": expected_assets,
            "discrepancy": discrepancy,
            "error_ratio": error_ratio,
            "is_valid": is_balanced,
            "risk_level": "LOW" if is_balanced else ("MEDIUM" if error_ratio < 0.05 else "CRITICAL")
        }
        return audit_verdict
\end{lstlisting}

---

\section{Đồ Án 3: Trợ Lý Chẩn Đoán Y Tế Hỗ Trợ Bác Sĩ (Clinical Diagnostic Copilot)}
\label{sec:workbook-clinical-copilot}

\subsection{Khởi Điểm Từ Con Số 0: Tiêu Chuẩn DICOM \& Phân Vùng Tổn Thương}
Dữ liệu y tế chẩn đoán hình ảnh (X-Quang, CT, MRI) được lưu trữ theo tiêu chuẩn quốc tế DICOM (Digital Imaging and Communications in Medicine). Không giống như ảnh JPEG 8-bit thông thường (chỉ có 256 mức xám), ảnh DICOM là dữ liệu 16-bit với hơn $65,536$ thang độ xám (đơn vị Hounsfield), cho phép phân biệt độ cản quang chính xác giữa mô mềm, dịch viêm và xương sườn.

Một trợ lý AI y tế không bao giờ được phép trực tiếp đưa ra kết luận khẳng định (\textit{"Bệnh nhân bị viêm phổi nặng"}), mà phải hoạt động theo mô hình hỗ trợ chẩn đoán: Khoanh vùng khu vực nghi ngờ tổn thương, trích xuất đặc trưng định lượng và gợi ý phác đồ y khoa chính thống từ cơ sở dữ liệu Bộ Y tế qua RAG.

\subsection{Mô Hình Toán Học: Hiệu Chuẩn Xác Suất Bằng Temperature Scaling}
Một lỗi phổ biến của các mạng nơ-ron sâu là hiện tượng tự tin thái quá (\textit{Overconfidence}): Mô hình dự đoán xác suất $99\%$ ngay cả khi hình ảnh mờ hoặc có nhiễu.
Để hiệu chuẩn xác suất dự đoán trước khi hiển thị cho bác sĩ, ta áp dụng kỹ thuật Temperature Scaling:
\begin{equation}
\hat{p}_i = \frac{\exp(z_i / T)}{\sum_{k=1}^K \exp(z_k / T)}
\end{equation}
với tham số nhiệt độ $T > 1$ được tối ưu hóa trên tập kiểm thử xác thực (\textit{Validation Set}) bằng cách cực tiểu hóa hàm mất mát Negative Log-Likelihood (NLL).

\subsection{Triển Khai Mã Nguồn Python: Bộ Hiệu Chuẩn Xác Suất \& Guardrails Y Tế}

\begin{lstlisting}[language=python]
import numpy as np

class ClinicalSafetyCalibrator:
    def __init__(self, temperature: float = 1.8):
        self.temperature = temperature
        # Cac cum tu cam doat quyen bac si
        self.prohibited_phrases = [
            "khang dinh 100%", "chac chan benh nhan bi", 
            "yeu cau uong thuoc ngay ma khong can y kien"
        ]

    def calibrate_probabilities(self, logits: np.ndarray) -> np.ndarray:
        """Hieu chuan lai phan phoi xac suat Softmax bang tham so Temperature."""
        scaled_logits = logits / self.temperature
        exp_logits = np.exp(scaled_logits - np.max(scaled_logits))
        return exp_logits / np.sum(exp_logits)

    def sanitize_clinical_recommendation(self, generated_text: str) -> dict:
        """Kiem tra ranh gioi an toan phap ly y te cua van ban tra loi."""
        violations = []
        for phrase in self.prohibited_phrases:
            if phrase in generated_text.lower():
                violations.append(phrase)
                
        is_safe = len(violations) == 0
        return {
            "is_safe": is_safe,
            "violations": violations,
            "sanitized_output": generated_text if is_safe else (
                "[CANH BAO AN TOAN Y TE]: Cau tra loi vi pham quy chuan tu van. "
                "Ket qua chi mang tinh chat tham khao, bac si can truc tiep tham kham."
            )
        }
\end{lstlisting}

---

\section{Đồ Án 4: Nền Tảng Gợi Ý Sản Phẩm Quy Mô Lớn (High-Throughput Two-Tower)}
\label{sec:workbook-two-tower-recsys}

\subsection{Khởi Điểm Từ Con Số 0: Kiến Trúc Two-Tower Cho Hệ Gợi Ý}
Trong một sàn thương mại điện tử có hàng triệu mặt hàng và hàng chục triệu người dùng, việc cho một mô hình học sâu phức tạp chấm điểm từng cặp (User, Item) theo thời gian thực là điều bất khả thi về mặt độ trễ ($O(M \times N)$ phép tính).

Kiến trúc Two-Tower giải quyết bài toán này bằng cách phân rã hệ thống thành hai nhánh độc lập:
\begin{enumerate}
  \item \textbf{User Tower (Nhánh Người Dùng)}: Nhận lịch sử duyệt web, vị trí địa lý, độ tuổi để ánh xạ thành véc-tơ người dùng $\mathbf{u} \in \mathbb{R}^d$.
  \item \textbf{Item Tower (Nhánh Mặt Hàng)}: Nhận hình ảnh, mô tả sản phẩm, danh mục để ánh xạ trước thành véc-tơ mặt hàng $\mathbf{v} \in \mathbb{R}^d$ và lưu trữ sẵn trong Vector DB.
\end{enumerate}
Khi người dùng truy cập trang chủ, hệ thống chỉ cần tính véc-tơ $\mathbf{u}$ một lần duy nhất, sau đó thực hiện tìm kiếm láng giềng gần nhất (ANN qua HNSW) trên tập véc-tơ $\mathbf{v}$ đã được chỉ mục từ trước, giảm độ trễ truy vấn xuống dưới $10\text{ms}$.

\subsection{Mô Hình Toán Học: Hàm Mất Mát InfoNCE (Contrastive Loss)}
Để huấn luyện hai nhánh mạng kéo gần các sản phẩm mà người dùng đã mua ($\mathbf{v}^+$) và đẩy xa các sản phẩm người dùng không quan tâm ($\mathbf{v}_j^-$), ta sử dụng hàm mất mát tương phản InfoNCE:
\begin{equation}
\mathcal{L}_{\text{InfoNCE}} = - \sum_{i=1}^B \log \frac{\exp\left(\frac{\mathbf{u}_i \cdot \mathbf{v}_i^+}{\tau}\right)}{\exp\left(\frac{\mathbf{u}_i \cdot \mathbf{v}_i^+}{\tau}\right) + \sum_{j \neq i} \exp\left(\frac{\mathbf{u}_i \cdot \mathbf{v}_j^+}{\tau}\right)}
\end{equation}
trong đó $\tau$ là siêu tham số nhiệt độ tương phản (\textit{Contrastive Temperature}), và các mẫu âm tính ($\mathbf{v}_j^+$ với $j \neq i$) được lấy ngay từ các sản phẩm của các người dùng khác trong cùng một mini-batch (\textit{In-Batch Negatives}), loại bỏ hoàn toàn chi phí lấy mẫu âm tính ngoại vi.

\subsection{Triển Khai Mã Nguồn Python: InfoNCE Loss \& Lập Chỉ Mục Vector}

\begin{lstlisting}[language=python]
import numpy as np

def compute_infonce_loss(user_embeddings: np.ndarray, 
                         item_embeddings: np.ndarray, 
                         temperature: float = 0.07) -> float:
    """Tinh toan ham mat mat In-Batch Negative InfoNCE Loss."""
    # Chuan hoa L2 cac vector embedding ve do dai 1
    u_norm = user_embeddings / (np.linalg.norm(user_embeddings, axis=1, keepdims=True) + 1e-9)
    v_norm = item_embeddings / (np.linalg.norm(item_embeddings, axis=1, keepdims=True) + 1e-9)
    
    # Ma tran tuong dong Cosine giua moi cap User va Item trong Batch: (B, B)
    similarity_matrix = (u_norm @ v_norm.T) / temperature
    
    # Mat na duong cheo: Cac phan tu u_i va v_i la cap duong tinh (Positive Pairs)
    batch_size = user_embeddings.shape[0]
    labels = np.arange(batch_size)
    
    # Ap dung Log-Sum-Exp Trick tinh Cross-Entropy Loss
    max_sim = np.max(similarity_matrix, axis=1, keepdims=True)
    exp_sim = np.exp(similarity_matrix - max_sim)
    sum_exp = np.sum(exp_sim, axis=1)
    
    # Lay phan tu duong cheo tuong ung voi cap positive
    pos_exp = exp_sim[np.arange(batch_size), labels]
    loss = - np.mean(np.log(pos_exp / sum_exp))
    return float(loss)
\end{lstlisting}

\subsection{Bộ Câu Hỏi Phản Biện Bảo Vệ Đồ Án}
\begin{enumerate}
  \item \textbf{Tại sao hệ gợi ý Two-Tower không thể so sánh tương tác chéo (\textit{Cross-Features}) giữa thuộc tính người dùng và sản phẩm?}
  \item \textbf{Hiện tượng phổ biến của In-Batch Negatives là \textit{Selection Bias} (các sản phẩm thịnh hành xuất hiện quá nhiều lần dưới dạng mẫu âm). Làm thế nào để hiệu chỉnh bằng xác suất lấy mẫu LogQ?}
  \item \textbf{Làm thế nào để phục vụ suy luận với thời gian trễ dưới 5ms khi tập sản phẩm mở rộng lên hơn 100 triệu phần tử?}
\end{enumerate}

---

\section{Đồ Án 5: Tường Lửa Proxy AI Tốc Độ Cực Cao Độ Trễ Thấp Dưới 5ms Cho Cụm Phục Vụ LLM (Ultra-Low Latency AI Firewall \& Prompt Guardrail Proxy)}
\label{sec:workbook-ai-firewall-proxy}

\subsection{Khởi Điểm Từ Con Số 0: Giao Thức HTTP Streaming \& Thách Thức Độ Trễ Guardrail}
Trong các ứng dụng trí tuệ nhân tạo hiện đại, trải nghiệm người dùng phụ thuộc sống còn vào chỉ số \term{TTFT (Time To First Token)} -- thời gian tính từ khi người dùng bấm gửi câu hỏi cho đến khi ký tự đầu tiên hiện ra trên màn hình. Để đạt được phản hồi mượt mà, các mô hình ngôn ngữ lớn (LLM) sử dụng cơ chế phát trực tuyến qua giao thức \term{Server-Sent Events (SSE)} trên nền HTTP/1.1 hoặc HTTP/2. Máy chủ gửi về client từng token một dưới dạng các khối bản tin liên tục:
\begin{center}
\code{data: \{"token": "Xin", "finish_reason": null\}}\\[2pt]
\code{data: \{"token": " chao", "finish_reason": null\}}
\end{center}

Tuy nhiên, một rào cản kỹ nghệ khổng lồ xuất hiện khi doanh nghiệp muốn bổ sung tầng bảo mật \term{Guardrails} để chống tấn công Prompt Injection, lộ lọt dữ liệu cá nhân (PII: số căn cước, mã thẻ ngân hàng), hoặc phát ngôn độc hại:
\begin{itemize}
  \item \textbf{Nghịch lý độ trễ của Guardrail LLM}: Nếu sử dụng một mô hình ngôn ngữ an toàn riêng biệt (ví dụ: Llama Guard 7B hoặc GPT-4o-mini) để đọc và phân loại prompt trước khi chuyển sang cụm phục vụ vLLM, thời gian phản hồi TTFT sẽ bị cộng thêm từ $250\text{ms}$ đến $600\text{ms}$. Đây là mức tăng độ trễ không thể chấp nhận được trong các hệ thống giao dịch tài chính hoặc trợ lý giọng nói trực tiếp.
  \item \textbf{Mục tiêu thiết kế}: Xây dựng một tầng \term{Reverse Guardrail Proxy} phi đồng bộ cắm ngay trước cụm vLLM. Proxy này kiểm duyệt hai chiều (cả prompt đầu vào và dòng token streaming đầu ra) với độ trễ cộng thêm \textbf{dưới $5\text{ms}$}, chặn đứng hoàn toàn các cuộc tấn công đã biết và ngắt luồng kết nối ngay tức khắc khi phát hiện nội dung vi phạm.
\end{itemize}

\subsection{Mô Hình Toán Học: Thuật Toán Aho-Corasick Đa Mẫu \& Phân Cách Siêu Phẳng Embedding}
Để đạt được tốc độ kiểm duyệt hàng triệu ký tự chỉ trong vài phần nghìn giây, ta kết hợp hai kỹ thuật toán học:

\paragraph{1. Thuật toán so khớp chuỗi hữu hạn Aho-Corasick}:
Cho một từ điển gồm $K$ mẫu nguy hại $\mathcal{P} = \{p_1, p_2, \dots, p_K\}$ với tổng độ dài các ký tự là $M = \sum |p_k|$. Một chuỗi văn bản đầu vào $\mathcal{T}$ có độ dài $N$ ký tự.
Nếu dùng thuật toán so khớp ngây thơ hoặc chạy $K$ biểu thức chính quy (Regex) độc lập, độ phức tạp tính toán là $\mathcal{O}(K \times N)$ -- chi phí sẽ bùng nổ khi tập mẫu luật tăng lên hàng chục ngàn từ.

Thuật toán Aho-Corasick ánh xạ toàn bộ tập từ điển $\mathcal{P}$ thành một cây tiền tố Trie kết hợp với hàm chuyển tiếp thất bại (\textit{Failure Link}):
\begin{itemize}
  \item \textbf{Hàm chuyển tiếp gốc $g(u, c)$}: Chuyển từ trạng thái $u$ sang trạng thái con khi gặp ký tự $c$.
  \item \textbf{Hàm thất bại $f(u)$}: Khi ký tự tiếp theo không có nhánh đi tiếp trong Trie, $f(u)$ trỏ về trạng thái tương ứng với hậu tố thực sự dài nhất của chuỗi ký tự dẫn tới $u$.
  \item \textbf{Hàm đầu ra $\text{out}(u)$}: Tập hợp tất cả các mẫu độc hại kết thúc tại trạng thái $u$ hoặc các hậu tố của $u$.
\end{itemize}
Nhờ cấu trúc đồ thị này, toàn bộ văn bản đầu vào $N$ ký tự được quét qua đúng một lần duy nhất theo thời gian tuyến tính:
\begin{equation}
\mathcal{T}_{\text{search}} = \mathcal{O}(N + Z)
\end{equation}
trong đó $Z$ là số lần xuất hiện của các mẫu vi phạm trong văn bản, hoàn toàn độc lập với kích thước từ điển $K$!

\paragraph{2. Phân loại ngữ nghĩa tuyến tính nhúng (Embedding Linear Hyperplane)}:
Đối với các đòn tấn công diễn đạt gián tiếp (Paraphrased Injection), ta sử dụng một mô hình nhúng siêu nhẹ (Tiny-Embedding 384 chiều) kết hợp phân loại nhị phân qua siêu phẳng Support Vector Machine:
\begin{equation}
\hat{y} = \sigma\left(\mathbf{w}^T \frac{\mathbf{e}(\mathcal{T})}{\|\mathbf{e}(\mathcal{T})\|_2} + b\right) \ge \theta_{\text{threshold}}
\end{equation}
Phép nhân vô hướng giữa véc-tơ 384 chiều và véc-tơ trọng số $\mathbf{w}$ trên tập lệnh SIMD CPU chỉ tiêu tốn chưa đầy $0.15\text{ms}$.

\subsection{Kiến Trúc Hệ Thống: Cửa Sổ Trượt Token \& Cắt Luồng Streaming (Streaming Cut-Off)}
Tường lửa hoạt động như một Reverse Proxy cắm giữa Client và Backend vLLM:
\begin{enumerate}
  \item \textbf{Kiểm duyệt đầu vào (Pre-Execution Ingress)}: Quét toàn bộ nội dung prompt qua máy Aho-Corasick và bộ phân loại tuyến tính. Nếu điểm nguy cơ vượt ngưỡng, Proxy trả về mã lỗi HTTP 403 Forbidden ngay lập tức trong $1.2\text{ms}$ mà không tiêu tốn một chu kỳ tính toán nào của GPU backend.
  \item \textbf{Kiểm duyệt dòng token đầu ra (Post-Execution Egress)}: Khi vLLM bắt đầu sinh dữ liệu, Proxy duy trì một \term{Sliding Window Ring Buffer} kích thước $W = 16$ tokens. Dòng token được truyền thông suốt tới người dùng để bảo đảm TTFT cực thấp. Ngay khi cửa sổ trượt phát hiện một mẫu vi phạm (ví dụ rò rỉ mã bí mật API Key hoặc thông tin y tế nhạy cảm), Proxy lập tức:
  \begin{itemize}
    \item Phát tín hiệu hủy luồng (\textit{Client Abort / Backend RST}) tới vLLM để giải phóng GPU KV-cache.
    \item Chèn thông báo thay thế \code{[REDACTED BY AI GUARDRAIL]} và ngắt dòng SSE với cờ \code{finish_reason: "content_filter"}.
  \end{itemize}
\end{enumerate}

\subsection{Triển Khai Mã Nguồn Python: Aho-Corasick Automaton \& Streaming Guardrail Proxy}

\begin{lstlisting}[language=python]
import time
from collections import deque
from dataclasses import dataclass, field
from typing import Dict, List, Optional, Tuple, Set

class AhoCorasickNode:
    """Nut trong cay tien to Aho-Corasick Automaton."""
    def __init__(self, depth: int = 0):
        self.depth = depth
        self.transitions: Dict[str, "AhoCorasickNode"] = {}
        self.fail: Optional["AhoCorasickNode"] = None
        self.outputs: List[str] = []

class AhoCorasickAutomaton:
    """Co che so khop da mau thoi gian tuyen tinh O(N + M)."""
    def __init__(self, keywords: List[str]):
        self.root = AhoCorasickNode(depth=0)
        for kw in keywords:
            self._insert(kw.lower())
        self._build_failure_links()

    def _insert(self, keyword: str) -> None:
        curr = self.root
        for char in keyword:
            if char not in curr.transitions:
                curr.transitions[char] = AhoCorasickNode(depth=curr.depth + 1)
            curr = curr.transitions[char]
        curr.outputs.append(keyword)

    def _build_failure_links(self) -> None:
        queue = deque()
        # Level 1 nodes have fail link to root
        for char, child in self.root.transitions.items():
            child.fail = self.root
            queue.append(child)

        # BFS construction of failure transitions
        while queue:
            curr = queue.popleft()
            for char, child in curr.transitions.items():
                queue.append(child)
                fail_node = curr.fail
                while fail_node and char not in fail_node.transitions:
                    fail_node = fail_node.fail
                child.fail = fail_node.transitions[char] if fail_node else self.root
                # Combine outputs with failure node
                child.outputs.extend(child.fail.outputs)

    def search(self, text: str) -> List[Tuple[int, str]]:
        """Tim tat ca cac mau doc hai trong doan van text."""
        results = []
        curr = self.root
        text_lower = text.lower()
        for idx, char in enumerate(text_lower):
            while curr and char not in curr.transitions:
                curr = curr.fail
            curr = curr.transitions[char] if curr else self.root
            if curr.outputs:
                for kw in curr.outputs:
                    results.append((idx - len(kw) + 1, kw))
        return results

@dataclass
class GuardrailVerdict:
    is_safe: bool
    latency_ms: float
    violations: List[str] = field(default_factory=list)
    action: str = "ALLOW"

class StreamingGuardrailProxy:
    """Tuong lua proxy AI do tre duoi 5ms cho cum vLLM."""
    def __init__(self, forbidden_signatures: List[str], window_size: int = 16):
        self.automaton = AhoCorasickAutomaton(forbidden_signatures)
        self.window_size = window_size

    def inspect_ingress_prompt(self, prompt: str) -> GuardrailVerdict:
        t0 = time.perf_counter()
        matches = self.automaton.search(prompt)
        elapsed_ms = (time.perf_counter() - t0) * 1000.0
        if matches:
            violations = list({m[1] for m in matches})
            return GuardrailVerdict(
                is_safe=False,
                latency_ms=elapsed_ms,
                violations=violations,
                action="BLOCK_403"
            )
        return GuardrailVerdict(is_safe=True, latency_ms=elapsed_ms, action="ALLOW")

    def inspect_streaming_chunk(self, ring_buffer: str, new_token: str) -> Tuple[bool, str, str]:
        """Kiem tra tung chunk token moi sinh ra tu engine."""
        combined_text = ring_buffer + new_token
        matches = self.automaton.search(combined_text)
        if matches:
            # Ngat luong va thay the bang token canh bao
            violation = matches[0][1]
            return False, violation, "[VIOLATION_INTERCEPTED]"
        return True, "", new_token
\end{lstlisting}

\subsection{Bộ Câu Hỏi Phản Biện Bảo Vệ Đồ Án}
\begin{enumerate}
  \item \textbf{Tại sao thuật toán Aho-Corasick có thể đảm bảo thời gian xử lý không đổi bất kể từ điển mở rộng từ 100 từ lên 100.000 từ?}
  \item \textbf{Nếu kẻ tấn công sử dụng kỹ thuật chèn ký tự vô hình (Zero-Width Space) hoặc ký tự đồng hình Unicode (Homoglyph) để vượt mặt Aho-Corasick, tầng Proxy cần bổ sung tiền xử lý gì?}
  \item \textbf{Làm thế nào để thiết kế bộ đệm trượt Ring Buffer cho dòng streaming SSE mà không làm tăng độ trễ cảm nhận của người dùng?}
\end{enumerate}

---

\section{Đồ Án 6: Nền Tảng Điều Phối Khôi Phục Thảm Họa Tự Động \& Dự Phòng Đa Vùng Cho Cụm AI (Multi-Region AI Disaster Recovery \& Autonomous Failover Orchestrator)}
\label{sec:workbook-multi-region-failover}

\subsection{Khởi Điểm Từ Con Số 0: Phân Vùng Mạng \& Hạ Tầng Phục Vụ AI Đa Vùng}
Trong môi trường sản xuất của các tập đoàn công nghệ lớn, các cụm máy chủ phục vụ mô hình ngôn ngữ lớn (vLLM / Triton clusters) không bao giờ được đặt tại duy nhất một trung tâm dữ liệu. Một sự cố mất điện diện rộng, đứt cáp quang biển quốc tế, hoặc lỗi cấu hình mạng BGP có thể khiến toàn bộ một vùng dữ liệu (\textit{Cloud Region} ví dụ: \code{us-east-1}) bị cô lập hoàn toàn khỏi Internet.

Hiện tượng một phần của mạng lưới bị mất liên lạc nhưng cả hai cụm máy chủ vẫn tiếp tục hoạt động độc lập được gọi là \term{Phân Vùng Mạng (Network Partition)}. Nếu không có cơ chế điều phối chặt chẽ, hệ thống rất dễ rơi vào trạng thái thảm họa \term{Split-Brain (Não Phân Liệt)}: Cả hai cụm máy chủ đều tự cho mình là cụm hoạt động duy nhất, cùng ghi đè trạng thái, gây mất đồng bộ dữ liệu và vi phạm các cam kết bảo mật.

Để giải quyết bài toán này, các Kỹ sư AI cần thiết kế một hệ thống điều phối \term{Autonomous Multi-Region Failover Orchestrator} có khả năng phát hiện sự cố nút mạng trong vài giây và chuyển hướng lưu lượng tải hoàn toàn tự động với cam kết:
\begin{itemize}
  \item \textbf{RTO (Recovery Time Objective) $\le 3\text{ giây}$}: Thời gian tối đa để hệ thống khôi phục hoàn toàn khả năng phục vụ truy vấn sau khi cụm chính sụp đổ.
  \item \textbf{RPO (Recovery Point Objective) $= 0$}: Không làm mất bất kỳ yêu cầu suy luận hợp lệ nào đang chờ trong hàng đợi.
\end{itemize}

\subsection{Mô Hình Toán Học: Thuật Toán Dự Báo Lỗi Phi-Accrual Failure Detector}
Cơ chế kiểm tra tình trạng sống còn (\textit{Health Check}) truyền thống dựa trên ngưỡng thời gian cố định (\textit{Fixed Timeout}, ví dụ: 5 giây không thấy nhịp tim thì coi như chết) có nhược điểm chí mạng:
\begin{itemize}
  \item Nếu đặt ngưỡng quá nhỏ (ví dụ 1 giây): Hiện tượng nghẽn mạng thoáng qua (\textit{Network Jitter}) sẽ kích hoạt chuyển vùng giả mạo (\textit{False Failover}), gây xáo trộn luồng kết nối hàng triệu người dùng.
  \item Nếu đặt ngưỡng quá lớn (ví dụ 10 giây): Thời gian phát hiện sự cố quá chậm, vi phạm chỉ tiêu SLA của doanh nghiệp.
\end{itemize}

Để giải quyết triệt để, ta áp dụng thuật toán dự báo lỗi thích ứng \term{$\Phi$-Accrual Failure Detector} (Hayashibara et al.):
Hệ thống duy trì một cửa sổ trượt $W = 100$ khoảng thời gian giữa các nhịp tim liên tiếp (\textit{Heartbeat Inter-arrival Times}) $\{x_1, x_2, \dots, x_W\}$.
Theo định lý giới hạn trung tâm, các biến ngẫu nhiên này có thể xấp xỉ bằng phân phối Gauss $\mathcal{N}(\mu, \sigma^2)$:
\begin{equation}
\mu = \frac{1}{W} \sum_{i=1}^W x_i, \quad \sigma^2 = \frac{1}{W-1} \sum_{i=1}^W (x_i - \mu)^2
\end{equation}

Tại thời điểm hiện tại $t$, nếu đã trôi qua một khoảng thời gian $\Delta = t - t_{\text{last}}$ kể từ lần cuối nhận được nhịp tim từ cụm máy chủ từ xa, xác suất để một nhịp tim bình thường xuất hiện muộn hơn thời điểm này là:
\begin{equation}
P_{\text{later}}(\Delta) = \frac{1}{\sigma \sqrt{2\pi}} \int_{\Delta}^{\infty} \exp\left(-\frac{(u - \mu)^2}{2\sigma^2}\right) du = 1 - \frac{1}{2} \left[ 1 + \text{erf}\left(\frac{\Delta - \mu}{\sigma \sqrt{2}}\right) \right]
\end{equation}

Mức độ nghi ngờ cụm máy chủ đã chết (\textit{Suspicion Level}) được định lượng trên thang đo Logarit:
\begin{equation}
\Phi = - \log_{10}\left( P_{\text{later}}(\Delta) \right)
\end{equation}

\textbf{Ý nghĩa định lượng của chỉ số $\Phi$}:
\begin{itemize}
  \item $\Phi = 1 \implies P_{\text{later}} = 0.1$ (khả năng trễ ngẫu nhiên là $10\%$).
  \item $\Phi = 3 \implies P_{\text{later}} = 10^{-3}$ (khả năng trễ ngẫu nhiên chỉ còn $0.1\%$).
  \item $\Phi \ge 8.0 \implies P_{\text{later}} \le 10^{-8}$ (sự chậm trễ này gần như chắc chắn $100\%$ là do trung tâm dữ liệu đã bị sập $\implies$ Kích hoạt quy trình chuyển vùng Failover tức thời!).
\end{itemize}

\paragraph{Phân tích độ sẵn sàng tổng thể hệ thống đa vùng}:
Xét hệ thống gồm $R$ vùng độc lập (ví dụ US-East, EU-Central, AP-Southeast), mỗi vùng có độ sẵn sàng độc lập $A_r = 99.9\%$ ($0.999$). Độ sẵn sàng song song toàn cục của toàn hệ thống là:
\begin{equation}
A_{\text{system}} = 1 - \prod_{r=1}^R (1 - A_r) = 1 - (1 - 0.999)^3 = 1 - (10^{-3})^3 = 1 - 10^{-9} = \mathbf{99.9999999\%} \quad (\text{9 số 9})!
\end{equation}

\subsection{Kiến Trúc Hệ Thống: Bầu Chọn Quorum \& Chuyển Hướng Lưu Lượng Anycast}
Hệ thống điều phối khôi phục thảm họa gồm 3 thành phần cốt lõi:
\begin{enumerate}
  \item \textbf{Cơ chế đồng thuận Quorum Trọng tài (Witness Arbiter)}: Để ngăn chặn Split-Brain khi đứt liên kết giữa Cụm 1 và Cụm 2, một thực thể trọng tài độc lập (\textit{Arbiter}) được đặt tại vùng thứ ba. Quyết định tiếp quản tải chỉ có hiệu lực khi đạt được sự đồng thuận của đa số tuyệt đối:
  \begin{equation}
  \text{Quorum} \ge \left\lfloor \frac{R}{2} \right\rfloor + 1
  \end{equation}
  \item \textbf{Chuyển hướng DNS \& BGP Anycast}: Khi kích hoạt chuyển vùng, Orchestrator cập nhật bản ghi DNS định tuyến độ trễ thấp (\textit{Latency-based DNS Routing}) qua Cloudflare/Route53 API và thu hồi tiền tố IP Anycast của cụm lỗi, chuyển hướng $100\%$ lưu lượng người dùng sang cụm dự phòng trong vòng $1.8\text{ giây}$.
\end{enumerate}

\subsection{Triển Khai Mã Nguồn Python: Phi-Accrual Detector \& Bộ Điều Phối Failover}

\begin{lstlisting}[language=python]
import math
import time
from collections import deque
from dataclasses import dataclass
from typing import Dict, List, Optional

class PhiAccrualFailureDetector:
    """Bo phat hien su co thich ung Phi-Accrual dua tren phan phoi Gauss."""
    def __init__(self, window_size: int = 100, min_std_dev: float = 0.05):
        self.window_size = window_size
        self.min_std_dev = min_std_dev
        self.intervals = deque(maxlen=window_size)
        self.last_heartbeat_time: Optional[float] = None

    def record_heartbeat(self, timestamp: float) -> None:
        """Ghi nhan nhip tim moi tu nut mang."""
        if self.last_heartbeat_time is not None:
            interval = timestamp - self.last_heartbeat_time
            if interval > 0:
                self.intervals.append(interval)
        self.last_heartbeat_time = timestamp

    def compute_phi(self, current_time: float) -> float:
        """Tinh toan muc do nghi ngo Phi tai thoi diem hien tai."""
        if self.last_heartbeat_time is None or len(self.intervals) < 3:
            return 0.0
        
        delta = current_time - self.last_heartbeat_time
        if delta <= 0:
            return 0.0

        mean = sum(self.intervals) / len(self.intervals)
        variance = sum((x - mean) ** 2 for x in self.intervals) / (len(self.intervals) - 1)
        std_dev = max(math.sqrt(variance), self.min_std_dev)

        # Tinh xac suat P_later qua ham sai so Gauss (erf)
        y = (delta - mean) / (std_dev * math.sqrt(2.0))
        # erf(y) xap xi tich phan Gauss
        p_later = 0.5 * math.erfc(y)
        
        # Tranh loi chia cho 0 hoac log(0)
        p_later = max(p_later, 1e-15)
        phi = - math.log10(p_later)
        return float(phi)

@dataclass
class RegionCluster:
    region_id: str
    is_primary: bool
    is_healthy: bool
    traffic_weight: float

class MultiRegionFailoverOrchestrator:
    """Dieu phoi chuyen vung tu dong cho cum phuc vu AI."""
    def __init__(self, primary_region: str, backup_region: str, phi_threshold: float = 8.0):
        self.clusters = {
            primary_region: RegionCluster(primary_region, is_primary=True, is_healthy=True, traffic_weight=1.0),
            backup_region: RegionCluster(backup_region, is_primary=False, is_healthy=True, traffic_weight=0.0)
        }
        self.primary_region = primary_region
        self.backup_region = backup_region
        self.detector = PhiAccrualFailureDetector()
        self.phi_threshold = phi_threshold
        self.failover_triggered = False

    def handle_heartbeat(self, timestamp: float) -> None:
        self.detector.record_heartbeat(timestamp)

    def evaluate_cluster_health(self, current_time: float) -> Tuple[bool, float, str]:
        """Danh gia suc khoe cum va tu dong chuyen vung neu can thiet."""
        phi = self.detector.compute_phi(current_time)
        if phi >= self.phi_threshold and not self.failover_triggered:
            # Phat hien su co tham hoa -> Chuyen giao luu luong
            self.failover_triggered = True
            self.clusters[self.primary_region].is_healthy = False
            self.clusters[self.primary_region].traffic_weight = 0.0
            
            self.clusters[self.backup_region].traffic_weight = 1.0
            return True, phi, f"CRITICAL_FAILOVER: Diverted 100% traffic to {self.backup_region}"
        
        return False, phi, "NORMAL"
\end{lstlisting}

\subsection{Bộ Câu Hỏi Phản Biện Bảo Vệ Đồ Án}
\begin{enumerate}
  \item \textbf{Tại sao thuật toán $\Phi$-Accrual ưu việt hơn cơ chế Heartbeat Timeout cố định khi đối mặt với hiện tượng mạng chập chờn (\textit{Network Jitter})?}
  \item \textbf{Trong bài toán phục vụ LLM, làm thế nào để đồng bộ bộ nhớ đệm tiền tố chung (\textit{Prefix KV-Cache}) giữa hai vùng địa lý cách nhau hàng ngàn kilomet mà không làm nghẽn băng thông?}
  \item \textbf{Cơ chế Quorum Arbiter hoạt động như thế nào để ngăn chặn hiện tượng não phân liệt (\textit{Split-Brain}) khi đường truyền giữa Cụm chính và Cụm dự phòng bị đứt gãy hoàn toàn?}
\end{enumerate}

---

\section{Đồ Án 7: Nền Tảng Hỗ Trợ Kỹ Thuật \& Tự Động Hóa Vận Hành IT / SRE Cho Doanh Nghiệp (Enterprise Autonomous IT Helpdesk \& SRE Copilot)}
\label{sec:project-7-sre-copilot}

\subsection{Bối Cảnh Doanh Nghiệp \& Thách Thức Vận Hành}
Tại các tập đoàn và doanh nghiệp công nghệ, bộ phận Hỗ trợ Kỹ thuật IT (IT Helpdesk) và Kỹ sư Vận hành Độ tin cậy Hệ thống (SRE / DevOps) luôn bị quá tải nghiêm trọng:
\begin{itemize}
  \item Hàng ngày tiếp nhận từ $500$ đến $2.000$ tickets hỗ trợ qua Slack, Microsoft Teams, và Jira Service Management.
  \item Khoảng $65\%$ số tickets là các tác vụ lặp đi lặp lại: Cấp lại mật khẩu VPN, phân quyền truy cập kho dữ liệu, kiểm tra tại sao container pod trên Kubernetes bị lỗi \code{CrashLoopBackOff} hoặc \code{OOMKilled}, hay kiểm tra độ trễ mạng chập chờn của chi nhánh văn phòng.
  \item Thời gian chờ phản hồi trung bình (First Response Time) kéo dài từ $2$ đến $6\text{ giờ}$, làm ngưng trệ năng suất làm việc của hàng nghìn nhân viên.
\end{itemize}

\subsection{Kiến Trúc Nền Tảng \& Vòng Lặp ReAct An Toàn}
Hệ thống kết hợp ba thành phần kiến trúc:
\begin{enumerate}
  \item \textbf{Phân Loại Mức Độ Nghiêm Trọng ITIL Tự Động}: Ánh xạ tức thì ticket vào 4 cấp độ: P1 (Sập dịch vụ toàn công ty), P2 (Ảnh hưởng phòng ban trọng yếu), P3 (Lỗi tính năng cục bộ), P4 (Yêu cầu tư vấn/thông tin).
  \item \textbf{Bộ Công Cụ Thao Tác Hệ Thống Chỉ Đọc (Read-Only Diagnostics)}: Agent được cấp quyền an toàn để gọi các API nội bộ: Kiểm tra trạng thái Kubernetes Pod, đọc 50 dòng log cuối cùng, đo độ trễ DNS, và tra cứu chính sách bảo mật nội bộ qua Hybrid RAG.
  \item \textbf{Điểm Ngắt Phê Duyệt Của Con Người (Human-in-the-Loop -- HITL)}: Đối với các thao tác thay đổi hạ tầng (Restart Pod, Flush Cache, Cấp quyền Admin), Agent bắt buộc phải dừng lại, xuất trình bản kế hoạch khắc phục sự cố và gửi nút xác nhận phê duyệt kèm chữ ký số tới Kỹ sư trực ca qua Slack Bot.
\end{enumerate}

\subsection{Hiện Thực Hóa Mã Nguồn Mẫu Bằng Python}
\begin{lstlisting}[language=Python]
from typing import Dict, Any, List, Optional, Tuple
import re

class EnterpriseSRECopilotEngine:
    """Dong co phan loai su co va tu dong dieu tra nguyen nhan goc (RCA)."""
    
    SEVERITY_KEYWORDS = {
        "P1_CRITICAL": ["down", "outage", "database crash", "sap he thong", "khong the thanh toan", "data breach"],
        "P2_HIGH": ["slow", "latency spike", "500 error", "oomkilled", "treo dich vu", "vpn failure"],
        "P3_MEDIUM": ["feature bug", "ui error", "cannot export pdf", "can cap quyen"],
        "P4_LOW": ["how to", "huong dan", "hoi dap", "general inquiry"]
    }

    @classmethod
    def classify_severity(cls, ticket_title: str, ticket_body: str) -> str:
        text = f"{ticket_title} {ticket_body}".lower()
        for sev, keywords in cls.SEVERITY_KEYWORDS.items():
            if any(kw in text for kw in keywords):
                return sev
        return "P3_MEDIUM"

    @classmethod
    def diagnose_pod_crash(cls, exit_code: int, termination_reason: str, log_snippet: str) -> Dict[str, Any]:
        """Chan doan nguyen nhan goc dua tren ma thoat va log Kubernetes."""
        root_cause = "UNKNOWN"
        remediation_action = "MANUAL_INVESTIGATION"
        requires_hitl = True

        if exit_code == 137 or "OOMKilled" in termination_reason:
            root_cause = "OUT_OF_MEMORY_LIMIT_EXCEEDED"
            remediation_action = "INCREASE_MEMORY_LIMIT"
            explanation = "Container bi Linux kernel OOM Killer tieu diet do vuot qua gioi han RAM duoc cap phat."
        elif exit_code == 1 or "CrashLoopBackOff" in termination_reason:
            if "database connection refused" in log_snippet.lower():
                root_cause = "UPSTREAM_DB_UNAVAILABLE"
                remediation_action = "CHECK_DB_NETWORK_CONNECTIVITY"
                explanation = "Ung dung khoi dong that bai do khong the ket noi toi may chu Database co so du lieu."
            else:
                root_cause = "APPLICATION_RUNTIME_EXCEPTION"
                remediation_action = "INSPECT_APP_STACK_TRACE"
                explanation = "Ung dung nem ra ngoai le nghiem trong trong qua trinh khoi tao."
        elif exit_code == 0:
            root_cause = "NORMAL_COMPLETION"
            remediation_action = "NO_ACTION_REQUIRED"
            requires_hitl = False
            explanation = "Tien trinh ket thuc thanh cong binh thuong."
        else:
            explanation = f"Container thoat voi ma loi khong xac dinh {exit_code}."

        return {
            "exit_code": exit_code,
            "root_cause": root_cause,
            "explanation": explanation,
            "recommended_action": remediation_action,
            "requires_human_approval": requires_hitl
        }

    @classmethod
    def generate_incident_postmortem(cls, incident_id: str, severity: str, root_cause: str, recovery_time_sec: float) -> str:
        """Sinh bao cao Post-Mortem su co tu dong dat chuan Google SRE."""
        return (
            f"=== BAO CAO PHAN TICH SU CO (POST-MORTEM) ===\n"
            f"Ma su co: {incident_id}\n"
            f"Muc do anh huong: {severity}\n"
            f"Nguyen nhan goc re (Root Cause): {root_cause}\n"
            f"Thoi gian khoi phuc (MTTR): {recovery_time_sec:.1f} giay\n"
            f"Tinh trang phe duyet: Da duoc chuyen vien SRE ky duyet\n"
            f"Hanh dong phong ngua: Cap nhat nguong canh bao Prometheus va bo sung gioi han tai nguyen."
        )
\end{lstlisting}

\subsection{Bộ Câu Hỏi Phản Biện Bảo Vệ Đồ Án}
\begin{enumerate}
  \item \textbf{Làm thế nào để bảo vệ hệ thống trước nguy cơ Prompt Injection nằm ẩn bên trong các đoạn log lỗi của container hoặc nội dung ticket người dùng gửi lên?}
  \item \textbf{Cơ chế nào ngăn chặn Agent rơi vào vòng lặp vô tận (Infinite ReAct Loop) khi thực hiện các lệnh chẩn đoán mạng bị timeout?}
  \item \textbf{Tại sao các hành động mang tính chất thay đổi trạng thái (\textit{Mutating Actions}) bắt buộc phải lưu vết chữ ký số mật mã học trên Audit Log?}
\end{enumerate}

---

\section{Đồ Án 8: Trợ Lý Giám Đốc Điều Hành Cá Nhân \& Tự Động Hóa Quy Trình (Executive AI Chief of Staff \& Workflow Automation)}
\label{sec:project-8-chief-of-staff}

\subsection{Bối Cảnh Cá Nhân / Doanh Nghiệp Nhỏ \& Bài Toán Đặt Ra}
Các nhà sáng lập startup, giám đốc điều hành và các chuyên gia cấp cao thường xuyên gặp phải tình trạng quá tải thông tin:
\begin{itemize}
  \item Nhận hàng trăm email mỗi ngày cần xử lý (từ báo giá đối tác, hóa đơn chi phí, lịch hẹn khách hàng đến thư quảng cáo rác).
  \item Lịch làm việc dày đặc dẫn tới tình trạng xung đột thời gian (\textit{Calendar Conflict}) giữa các cuộc hẹn quan trọng.
  \item Quản lý chứng từ chi tiêu và hóa đơn tài chính rời rạc, làm chậm tiến độ quyết toán thuế hàng tháng.
\end{itemize}

\subsection{Kiến Trúc Giải Pháp \& Ba Trụ Cột Tự Động Hóa}
\begin{enumerate}
  \item \textbf{Hộp Thư Thông Minh (Inbox Zero AI Engine)}: Trích xuất ý định (\textit{Intent Recognition}), nhận diện mức độ khẩn cấp, tóm tắt 1 dòng cốt lõi và tự động dự thảo thư phản hồi (\textit{Draft Response}) bám sát văn phong của người dùng.
  \item \textbf{Bộ Điều Phối Lịch Trình (Smart Calendar Conflict Resolver)}: Sử dụng thuật toán giao thoa khoảng thời gian đóng mở $[s_1, e_1] \cap [s_2, e_2] \neq \emptyset$ để phát hiện trùng lịch và tự động đề xuất 3 khung thời gian trống khả dụng chung.
  \item \textbf{Trích Xuất Hóa Đơn Tài Chính Đa Phương Thức (Vision Invoice Parser)}: Đọc dữ liệu từ hình ảnh hoặc PDF hóa đơn VAT, trích xuất mã số thuế, nhà cung cấp, ngày lập và số tiền thanh toán, sau đó kiểm tra tính hợp lệ của công thức thuế VAT $10\%$.
\end{enumerate}

\subsection{Hiện Thực Hóa Mã Nguồn Mẫu Bằng Python}
\begin{lstlisting}[language=Python]
from typing import Dict, Any, List, Optional, Tuple

class ExecutiveChiefOfStaffEngine:
    """Dong co ho tro giam doc dieu hanh xu ly email, lich trinh va hoa don."""
    
    @staticmethod
    def triage_email(sender: str, subject: str, body: str) -> Dict[str, Any]:
        text = f"{subject} {body}".lower()
        
        # 1. Phan loai danh muc
        if any(w in text for w in ["invoice", "hoa don", "thanh toan", "receipt", "billing"]):
            category = "FINANCE_EXPENSE"
            priority = "HIGH"
        elif any(w in text for w in ["gap mat", "hop", "meeting", "calendar", "call"]):
            category = "CALENDAR_REQUEST"
            priority = "HIGH"
        elif any(w in text for w in ["khuyen mai", "sale", "discount", "newsletter", "unsubscribe"]):
            category = "MARKETING_SPAM"
            priority = "LOW"
        else:
            category = "GENERAL_COMMUNICATION"
            priority = "MEDIUM"

        # 2. Sinh tom tat va du thao hanh dong
        is_action_required = priority in ["HIGH", "MEDIUM"]
        return {
            "sender": sender,
            "subject": subject,
            "category": category,
            "priority": priority,
            "requires_action": is_action_required
        }

    @staticmethod
    def detect_calendar_conflict(existing_events: List[Tuple[float, float]], new_start: float, new_end: float) -> Tuple[bool, Optional[Tuple[float, float]]]:
        """Kiem tra xem su kien moi [new_start, new_end] co bi trung voi lich cu hay khong."""
        if new_start >= new_end:
            raise ValueError("Thoi gian bat dau phai nho hon thoi gian ket thuc")
            
        for (start, end) in existing_events:
            # Hai khoang [s1, e1] va [s2, e2] giao nhau khi: max(s1, s2) < min(e1, e2)
            if max(start, new_start) < min(end, new_end):
                return True, (start, end)
        return False, None

    @staticmethod
    def parse_and_validate_invoice(subtotal: float, vat_rate: float, total_billed: float, supplier_tax_id: str) -> Dict[str, Any]:
        """Kiem tra tinh hop le cua hoa don VAT."""
        if subtotal < 0 or vat_rate < 0 or total_billed < 0:
            raise ValueError("So tien khong the am")
            
        expected_vat = subtotal * vat_rate
        expected_total = subtotal + expected_vat
        discrepancy = abs(total_billed - expected_total)
        is_valid = discrepancy < 0.01  # Cho phep sai so lam tron 1 xu
        
        # Kiem tra dinh dang Ma So Thue Viet Nam (10 hoac 13 chu so)
        is_tax_id_valid = bool(re.match(r"^\d{10}(-\d{3})?$", supplier_tax_id))

        return {
            "subtotal": subtotal,
            "vat_amount": expected_vat,
            "expected_total": expected_total,
            "discrepancy": round(discrepancy, 4),
            "is_math_valid": is_valid,
            "is_tax_id_valid": is_tax_id_valid,
            "status": "APPROVED" if (is_valid and is_tax_id_valid) else "FLAGGED_FOR_REVIEW"
        }
\end{lstlisting}

\subsection{Bộ Câu Hỏi Phản Biện Bảo Vệ Đồ Án}
\begin{enumerate}
  \item \textbf{Làm thế nào để đảm bảo tính riêng tư (Privacy) tuyệt đối cho email và dữ liệu tài chính của giám đốc khi sử dụng mô hình LLM bên ngoài qua API đám mây?}
  \item \textbf{Nếu hai đối tác cùng gửi yêu cầu hẹn gặp vào cùng một khung giờ trống, thuật toán thương lượng lịch trình (Calendar Negotiation) sẽ giải quyết ưu tiên như thế nào?}
  \item \textbf{Trong bài toán trích xuất hóa đơn bằng Vision-LLM, làm thế nào để phát hiện và ngăn chặn các hóa đơn giả mạo được tạo bởi các công cụ AI sinh ảnh?}
\end{enumerate}


